\documentclass[preprint,12pt,authoryear]{elsarticle}

\usepackage{amsmath,amssymb,amsthm}
\usepackage{booktabs}
\usepackage{multirow}
\usepackage{xcolor}
\usepackage{graphicx}
\usepackage{hyperref}

\newtheorem{proposition}{Proposition}
\newtheorem{remark}{Remark}
\newcommand{\placeholder}[1]{\textcolor{red}{[#1]}}
\newcommand{\norm}[1]{\left\lVert #1 \right\rVert}

\journal{Transportation Research Part E: Logistics and Transportation Review}

\begin{document}

\begin{frontmatter}

\title{Deep Reinforcement Learning with Sparse Exact Feedback for the Carrier--Vehicle Traveling Salesman Problem}

\author{\placeholder{Author names, affiliations and corresponding author}}

\begin{abstract}
In the carrier--vehicle traveling salesman problem (CVTSP), a carrier launches and recovers a drone at continuous points while the drone visits a set of targets. Evaluating one visiting order requires solving a second-order cone program (SOCP), so training a deep reinforcement learning (DRL) policy on exact rewards consumes hundreds of millions of solver calls. We show that one optimal SOCP solution yields a certified upper bound on the cost of every visiting order, because fixed take-off and landing points remain feasible when the order changes. We use this bound as the main training signal: it supplies dense rewards, guides solver-free local search and screens candidate expert routes, while exact solves are reserved for a few routes per instance. Compared with POMO trained on exact rewards for every route, the proposed method uses 15 times fewer training instances and 78 times fewer solver calls at the same test time. On 529 literature benchmark instances its mean per-instance gap to that baseline is +0.21\%. The gap is +2.60\% on 500 synthetic instances without augmentation and +0.72\% on instances with 150 to 500 targets with eightfold augmentation. Long-budget metaheuristics obtain better solutions when given ten times the computing time. The method is therefore most useful when training budgets are limited and solutions are needed quickly.
\end{abstract}

\begin{keyword}
Carrier--vehicle routing \sep Drone logistics \sep Deep reinforcement learning \sep Second-order cone programming \sep Certified upper bound
\end{keyword}

\end{frontmatter}

\section{Introduction}
\label{sec:intro}

A drone with a short flight range can serve targets far beyond that range if a long-range platform carries it. A ship can launch an unmanned aerial vehicle (UAV) towards an offshore installation, keep moving, and recover the UAV further along its course. The carrier--vehicle traveling salesman problem (CVTSP), introduced by \citet{garone2011}, models this operation. A carrier leaves a depot, releases a drone before each target, recovers it after the visit and returns to the depot once all targets have been served. The planner chooses the visiting order and the take-off and landing point of each sortie in order to minimise mission completion time.

The problem combines two kinds of decisions. The visiting order is a permutation, and the sortie points lie in the continuous plane. Once the order is fixed, the remaining problem is convex and can be written as an SOCP \citep{lobo1998}. The joint problem is a mixed-integer SOCP. \citet{gambella2018} and \citet{erdogan2021} solved it exactly for small instances, and \citet{vieira2026} report optimal solutions for all benchmark instances with up to 25 customers and for one instance with 35 customers. Larger instances rely on heuristics such as the two-stage iterative heuristic of \citet{li2025}. \citet{ren2026} extend the model to a vessel that carries several UAVs, which shows that the setting remains an active research problem in maritime logistics.

Deep reinforcement learning has become a practical alternative for routing problems when solutions must be produced quickly. A trained policy produces a route in a single forward pass, and methods such as POMO \citep{kwon2020} reach small gaps on the traveling salesman problem. These methods rest on a property that the CVTSP lacks: the cost of a route is a sum of known edge lengths. In the CVTSP, the cost of a route is the optimal value of an SOCP. Our reference baseline, POMO trained on an exact solver reward for every sampled route, made about $3 \times 10^{8}$ solver calls over 200 epochs. That cost must be paid again whenever the instance distribution or the operating constraints change. For a maritime or humanitarian operator, a change of vessel, drone model or service region is a change of distribution, so training cost is a recurring operational expense and not a one-off research cost.

This paper asks whether most of those calls are necessary. The answer rests on one structural observation. In the fixed-order SOCP, the constraints on a sortie concern one target only, and the visiting order enters only through the carrier legs between sorties. If the sortie points from one optimal solution are kept and the order is changed, every sortie remains feasible and only the carrier legs change. The resulting value is an upper bound on the optimal cost of the new order, and it can be computed by adding up entries of a table. One exact solve thus says something about every route of the instance, including the routes that will never be solved.

In the literature we searched (OpenAlex, arXiv, Crossref and web sources, last searched 4 November 2026), we found no DRL method for routing problems with a continuous-optimisation route cost that reuses one exact primal solution as a certified upper bound on all visiting orders, for both training rewards and expert search. The closest work is \citet{li2025}, whose heuristic also derives an edge matrix from solved subproblems. Their matrix updates traversed arcs only, keeps estimates elsewhere and carries no bound. We return to this comparison in Section~\ref{sec:gap}.

The paper makes three contributions.
\begin{enumerate}
\item We state and prove that one optimal fixed-order solution yields a table whose route sums bound the exact cost of every visiting order from above and match it at the solved order. The proof is short because it uses the structure of the CVTSP directly. The contribution lies in how this bound is put to work.
\item We build a training method around the table. Dense rewards, local search and expert screening use the table without calling a solver, and exact solves are reserved for the anchor route, for one policy route and for expert verification. The method uses 2.89 solver calls per training instance on average.
\item We compare the method with POMO trained on exact rewards for every route, and we report the results by instance size. The method uses 15 times fewer training instances and 78 times fewer solver calls, with no change in test time. It remains behind the baseline on synthetic instances, and we quantify where and by how much.
\end{enumerate}

The practical result is a change in what a learned CVTSP solver costs to obtain. Training on a single workstation finishes in less than a day, compared with more than a week for the baseline on a larger worker pool. The paper also reports where this comes at a price: the method does not match the baseline on instances from its own training distribution, and metaheuristics with long time budgets obtain better solutions. We report both results alongside the savings, because a planner choosing a method needs all three.

Section~\ref{sec:related} reviews related work. Section~\ref{sec:problem} defines the problem and Section~\ref{sec:method} presents the method. Section~\ref{sec:experiments} reports the experiments, Section~\ref{sec:discussion} discusses practical implications and Section~\ref{sec:conclusion} concludes.

\section{Related work}
\label{sec:related}

Three strands of work bear on this paper: solution methods for the CVTSP, DRL for routing, and learning under expensive objectives.

\subsection{The CVTSP and its extensions}

\citet{garone2011} introduced the problem and proposed a two-step approach that first fixes an order and then optimises the sortie points. \citet{garone2014} extended it to carriers that visit several targets per sortie. \citet{gambella2018} gave the first exact mixed-integer SOCP approach, solving instances with up to 15 targets. \citet{erdogan2021} tightened the formulation and added an iterated local search, and \citet{poikonen2020} studied the closely related mothership and drone routing problem with instances of up to 200 targets. \citet{fahradyan2020} considered several carrier--vehicle pairs, and \citet{liy2024} studied a joint vessel--UAV routing variant with a tabu search heuristic.

Recent work has moved the exact and heuristic frontiers. \citet{li2025} developed exact methods and a two-stage iterative heuristic, and their results supply the reference values used in this paper. \citet{vieira2026} propose a linearisation-based refinement algorithm that solves all benchmark instances with up to 25 customers. \citet{ren2026} address a vessel with multiple UAVs through a mixed-integer SOCP and an adaptive large neighbourhood search. CVTSP research is one branch of a wider literature on drone-assisted routing \citep{murray2015,agatz2018,otto2018,macrina2020,chung2020}, and transport journals have recently paid attention to operational uncertainty in such systems \citep{zhao2022}. All of the CVTSP methods above solve each instance from scratch. They spend their computing time per instance, whereas a trained policy pays once during training.

Two observations from this literature shape our design. The first is that most of the heuristics we reviewed rest on the fixed-order problem: a candidate order has to be solved before its true cost is known. The second is that the order and the sortie points are strongly coupled, since a good order depends on where the carrier will be, and the carrier's position depends on the order. Two-stage methods handle this coupling by alternating between the two levels. Our table can be read as one such alternation step, frozen and reused across a whole search.

\subsection{Deep reinforcement learning for routing}

Learning to construct routes began with pointer networks \citep{vinyals2015} and was extended to vehicle routing by \citet{nazari2018}. The attention model of \citet{kool2019} established the encoder--decoder architecture that most later work adopts. POMO \citep{kwon2020} exploits the multiple symmetric starts of a route, and Sym-NCO \citep{kim2022} and efficient active search \citep{hottung2022} add symmetry training and instance-level adaptation. Improvement-based policies \citep{wu2022} and heavy-decoder models \citep{luo2023} target better quality and larger sizes. Surveys by \citet{bengio2021} and \citet{mazyavkina2021} cover the field, and \citet{wuy2026} show that this line of work is reaching transport journals through clustered vehicle routing. Across this body of work, the reward signal is rarely a concern, since an additional route evaluation costs a few floating-point operations.

Most of these solvers evaluate a route in closed form, as a sum of distances with a feasibility check. Work on routing in continuous space comes closer to our setting. \citet{fan2025} address the close-enough TSP, in which a tour only needs to reach a neighbourhood of each target. They discretise candidate waypoints, so that the reward remains a tour length and the training loop avoids solving a continuous subproblem. \citet{liu2026} apply reinforcement learning to eVTOL-and-drone routing in continuous space, \citet{yan2026} study truck-based mothership delivery to moving targets with multi-agent reinforcement learning, and \citet{lee2023} apply reinforcement learning to mothership routing. \placeholder{Confirm the methods of Liu, Yan and Lee from their full texts; the wording here rests on titles and search summaries.} Within the sources we searched, we found no learning method stated for the CVTSP as defined by \citet{garone2011}. That finding is bounded by our search and does not extend to mothership routing in general. DRL has clearly reached related problems, so we do not argue that this class of problem is untouched. Our argument concerns the cost of training when each reward requires a solver.

\subsection{Expensive objectives and self-improvement}

Expensive objectives have received direct attention in neural combinatorial optimisation. \citet{kim2024} propose symmetric replay training, which reuses solutions that have already been evaluated through their symmetries and thereby saves objective evaluations. It offers no estimate for a route that was never evaluated. \citet{sun2026} train a neural network offline to estimate drone flight times from trajectory-optimisation data and use it inside a heuristic for truck and drone delivery. SurCo \citep{ferber2023} learns linear surrogate costs for nonlinear combinatorial problems. These methods learn or reuse a cost signal, and none of those we read provides a guarantee on routes outside the evaluated set. A learned surrogate can be accurate on average, but it can also be wrong in a direction that is hard to detect, and its training set has to be produced in advance. A bound derived from the solution of the instance at hand needs no training data and has a known direction of error, which is what makes it usable as an acceptance test.

Self-improvement offers a second route to lower training cost, in which a policy trains on labels that it produces itself. \citet{luo2023} train by supervised learning on small instances, \citet{pirnay2024} and \citet{luo2024} generate pseudo-labels from the current model, and \citet{luttmann2025} extend the idea to multi-action settings. On the problems these papers study, comparing two candidate solutions is cheap, so labels can be accepted or rejected by direct evaluation. In the CVTSP, the same test costs one SOCP per candidate. Our method keeps the self-improvement loop and uses the bound to decide which candidates deserve that cost. This changes the economics of the loop. Generating a hundred candidates is cheap, but verifying a hundred candidates is not, so the screening step becomes the part of the method that determines its budget.

\subsection{Research gap}
\label{sec:gap}

Table~\ref{tab:gap} compares the closest work on five properties. Among the sources we read, none combines a continuous route-cost subproblem, a training signal for routes that were never solved, and a guarantee on that signal. The contrast with \citet{li2025} is the sharpest. Their approximate matrix is initialised in closed form and, after a fixed-order problem is solved, updated only on traversed arcs; the other entries keep their prior values. The matrix guides a heuristic and is regenerated for diversification. Our table is built from one primal solution in a single step, assigns a valid value to every arc and bounds the cost of every order. Fixing part of a solution and reoptimising the rest is an established pattern in routing, so we do not claim Proposition~\ref{prop:bound} as a new idea in itself. Our claim concerns its use as the main training signal for a DRL policy.

\begin{table}[t]
\centering
\caption{Closest work compared on properties relevant to this paper. Liu et al. (2026), Yan et al. (2026) and Lee et al. (2023) are omitted because their full texts have not been read.}
\label{tab:gap}
\small
\resizebox{\textwidth}{!}{\begin{tabular}{lccccc}
\toprule
Work & \begin{tabular}[c]{@{}c@{}}Continuous\\ subproblem\end{tabular} & \begin{tabular}[c]{@{}c@{}}Exact feedback\\ in training\end{tabular} & \begin{tabular}[c]{@{}c@{}}Value for\\ unsolved routes\end{tabular} & \begin{tabular}[c]{@{}c@{}}Bound on\\ all orders\end{tabular} & \begin{tabular}[c]{@{}c@{}}DRL builds\\ the route\end{tabular} \\
\midrule
\citet{li2025} & Yes & No & Estimate & No & No \\
\citet{vieira2026} & Yes & No & No & No & No \\
\citet{ren2026} & Yes & No & No & No & No \\
\citet{kim2024} & No & No & No & No & Yes \\
\citet{sun2026} & Yes & Offline data & Estimate & No & No \\
\citet{fan2025} & Discretised & No & Closed form & No & Yes \\
\citet{pirnay2024}, \citet{luo2024} & No & No & Closed form & No & Yes \\
This paper & Yes & Yes, sparse & Certified & Yes & Yes \\
\bottomrule
\end{tabular}}
\end{table}

\section{Problem formulation}
\label{sec:problem}

Let $d \in \mathbb{R}^{2}$ be the depot and $p_{1},\dots,p_{J} \in \mathbb{R}^{2}$ the targets. The carrier moves at speed $v_{c}$ and the drone at speed $v_{d} > v_{c}$. The drone can fly for at most $E$ time units per sortie. Each target must be visited in exactly one sortie. For target $i$, the drone leaves the carrier at take-off point $s_{i}$, flies to $p_{i}$ and returns to the carrier at landing point $l_{i}$. The carrier travels from $s_{i}$ to $l_{i}$ during the sortie.

\subsection{The fixed-order problem}

Let $\pi = (\pi_{1},\dots,\pi_{J})$ be a visiting order, with the conventions that $l_{\pi_{0}}$ and $s_{\pi_{J+1}}$ denote the depot. Let $t^{1}_{i}$ and $t^{2}_{i}$ be the outbound and inbound flight times, $\tau_{i}$ the sortie duration, and $T_{k}$ the carrier travel time on the leg between consecutive sorties. The carrier--vehicle problem for order $\pi$ is
\begin{align}
f(\pi) = \min \quad & \sum_{i=1}^{J} \tau_{i} + \sum_{k=0}^{J} T_{k} \label{eq:obj} \\
\text{s.t.} \quad & \norm{s_{i} - p_{i}} \le v_{d}\, t^{1}_{i}, && i = 1,\dots,J, \label{eq:out}\\
& \norm{l_{i} - p_{i}} \le v_{d}\, t^{2}_{i}, && i = 1,\dots,J, \label{eq:in}\\
& \norm{l_{i} - s_{i}} \le v_{c}\, \tau_{i}, && i = 1,\dots,J, \label{eq:sync}\\
& \tau_{i} = t^{1}_{i} + t^{2}_{i}, \quad \tau_{i} \le E, && i = 1,\dots,J, \label{eq:end}\\
& \norm{s_{\pi_{k+1}} - l_{\pi_{k}}} \le v_{c}\, T_{k}, && k = 0,\dots,J, \label{eq:leg}\\
& t^{1}_{i},\, t^{2}_{i},\, \tau_{i},\, T_{k} \ge 0 . \label{eq:nonneg}
\end{align}
Every constraint is a second-order cone or a linear constraint, so the problem is an SOCP \citep{lobo1998}. The CVTSP asks for $\min_{\pi} f(\pi)$. Mission completion time is the objective in \eqref{eq:obj}, because the carrier is always either on a leg or within a sortie. Table~\ref{tab:notation} lists the main notation.

\begin{table}[t]
\centering
\caption{Main notation.}
\label{tab:notation}
\small
\begin{tabular}{ll}
\toprule
Symbol & Meaning \\
\midrule
$J$ & Number of targets \\
$d$, $p_{i}$ & Depot and target locations \\
$v_{c}$, $v_{d}$ & Carrier and drone speeds \\
$E$ & Maximum sortie duration \\
$\pi$ & Visiting order, with $\pi_{0} = \pi_{J+1} = 0$ denoting the depot \\
$s_{i}$, $l_{i}$ & Take-off and landing points of the sortie for target $i$ \\
$t^{1}_{i}$, $t^{2}_{i}$, $\tau_{i}$ & Outbound time, inbound time and duration of a sortie \\
$T_{k}$ & Carrier travel time on leg $k$ \\
$f(\pi)$ & Optimal value of the fixed-order problem for $\pi$ \\
$q$ & Anchor route, solved exactly \\
$C^{q}$, $F^{q}$ & Directed table and total sortie time from the solution for $q$ \\
$U_{q}(\pi)$ & Certified upper bound on $f(\pi)$ \\
\bottomrule
\end{tabular}
\end{table}

\subsection{Cost structure}

Two features of this formulation drive the method. First, constraints \eqref{eq:out}--\eqref{eq:end} involve the variables of a single target and make no reference to $\pi$. Second, the order appears only in \eqref{eq:leg}, which links the landing point of one target to the take-off point of the next. Any solution that fixes all sortie variables therefore induces a closed-form cost for every order, since the leg times are the only remaining quantities. Evaluating a fixed-order problem with $J$ targets takes tens of milliseconds at moderate sizes in our implementation, about 0.03\,s per call at $J \approx 60$, while evaluating a table sum takes microseconds. The number of orders that can be scored in one solve time differs by several orders of magnitude, which is the leverage the method relies on. This is the structure that the next section exploits. The policy in Section~\ref{sec:method} selects $\pi$, and the SOCP determines the sortie points, so the method follows the natural decomposition into a discrete upper level and a continuous lower level. Only the upper level is learned.

\section{Methodology}
\label{sec:method}

\subsection{Overview}

Table~\ref{alg:train} summarises one training step for a batch of instances. A frozen reference policy samples 200 candidate routes per instance. A small edge predictor selects one of them as the anchor, and the anchor is solved exactly. The solution yields a table $C$ that scores every candidate. Local search on the table produces an expert route, which is verified exactly before it can serve as an imitation label. The learning policy then receives two updates: a reinforcement step on rewards drawn mainly from the table, and an imitation step on the verified experts. The number of exact solves per instance is independent of the number of POMO routes. This independence is the source of the savings. In POMO with an exact reward, a batch of $B$ instances with $P$ routes costs $BP$ solves. Here it costs about $3B$, plus a refresh term for large instances, regardless of $P$.

\begin{table}[t]
\caption{One training step for a batch of instances. Exact solves are marked.}
\label{alg:train}
\small
\hrule\smallskip
\noindent For each instance in the batch:
\begin{enumerate}
\setlength{\itemsep}{0pt}
\item Sample 200 routes from the frozen reference policy.
\item Rank them with the edge predictor and choose the anchor $q$.
\item \textbf{Exact solve:} solve the SOCP for $q$ and build the table $C^{q}$.
\item Score all sampled routes with $U_{q}$ and select 30 starts.
\item Run 2-opt on $C^{q}$ from each start, and LKH from a share of the starts when $J \ge 60$.
\item \textbf{Exact solve:} if the best route has $U_{q}$ below $f(q)$, solve it. Otherwise, with the near-miss probability, solve the closest route.
\item \textbf{Exact solves:} for $J \ge 60$ after epoch 60, run up to two refresh rounds of the search, with one solve per round.
\end{enumerate}
\noindent For the batch as a whole:
\begin{enumerate}
\setlength{\itemsep}{0pt}
\setcounter{enumi}{7}
\item Decode the POMO routes of the learning policy and score them with the tables.
\item \textbf{Exact solve:} solve the best-scoring route of each instance.
\item Update the policy by REINFORCE on the hybrid reward.
\item Update the policy by cross-entropy on the verified expert routes.
\item Update the edge predictor on the tables.
\end{enumerate}
\hrule
\end{table}

\subsection{A certified upper bound on all visiting orders}

Let $q$ be a visiting order and let $(s^{*},l^{*},\tau^{*},T^{*})$ be an optimal solution of the fixed-order problem for $q$, with each sortie variable indexed by target. Write $F^{q} = \sum_{i} \tau^{*}_{i}$ and define a directed table on the nodes $\{0,1,\dots,J\}$, with node $0$ the depot:
\begin{equation}
C^{q}_{0j} = \frac{\norm{s^{*}_{j} - d}}{v_{c}}, \qquad
C^{q}_{i0} = \frac{\norm{d - l^{*}_{i}}}{v_{c}}, \qquad
C^{q}_{ij} = \frac{\norm{s^{*}_{j} - l^{*}_{i}}}{v_{c}} .
\end{equation}
For any order $\pi$, define the route value
\begin{equation}
U_{q}(\pi) = F^{q} + \sum_{k=0}^{J} C^{q}_{\pi_{k}\pi_{k+1}}, \qquad \pi_{0} = \pi_{J+1} = 0 .
\end{equation}

\begin{proposition}
\label{prop:bound}
For every visiting order $\pi$, $f(\pi) \le U_{q}(\pi)$. Moreover, $f(q) = U_{q}(q)$.
\end{proposition}

\begin{proof}
Fix $\pi$. Keep $s^{*}_{i}$, $l^{*}_{i}$, $t^{1*}_{i}$, $t^{2*}_{i}$ and $\tau^{*}_{i}$ for each target and set $T_{k} = C^{q}_{\pi_{k}\pi_{k+1}}$. Constraints \eqref{eq:out}--\eqref{eq:end} involve only these sortie variables and do not depend on the order, so they hold because they hold for $q$. Constraint \eqref{eq:leg} holds with equality by construction, and non-negativity is immediate. The point is thus feasible for the fixed-order problem of $\pi$, with objective value $U_{q}(\pi)$. Since $f(\pi)$ is the minimum over feasible points, $f(\pi) \le U_{q}(\pi)$. For $\pi = q$, each $T^{*}_{k}$ is at its lower bound at an optimal solution because the objective strictly increases in $T_{k}$, which gives $U_{q}(q) = f(q)$.
\end{proof}

Three conditions apply in practice. The anchor must be solved to optimality; a merely feasible solution would still give a valid upper bound, but the equality at $q$ could fail. In the implementation, every anchor solution passes an audit of all constraints with a tolerance of $10^{-5}$ before its table is used, and the bound therefore holds up to that tolerance. The training code also checks the inequality whenever an expert, near-miss or refresh route is solved exactly, and stops on a violation. The final run completed without a stop.

\begin{remark}
Several tables can be combined only by minimising over complete routes, $\min_{q} U_{q}(\pi)$. An entry-wise minimum of tables would splice together sortie points from different solutions and loses the guarantee, because the implied take-off and landing point of one target would differ between its incoming and outgoing arcs.
\end{remark}

The table is asymmetric, so local search must account for the reversal of internal arcs. With prefix sums of the differences between forward and reverse arc costs, a 2-opt move is scored in constant time, without any solver call.

\subsection{Policy network}

The learning policy is an attention encoder--decoder of the POMO type \citep{kwon2020}, with 4 encoder layers, embedding dimension 256, 16 attention heads and feed-forward dimension 512. Node features are the depot-relative coordinates scaled by instance size, the distance to the depot, a depot indicator, the speed ratio $v_{c}/v_{d}$ and the drone reach $v_{d}E$ relative to the instance scale. The speed and endurance features let one network serve instances with different operating parameters. The policy decodes one route per POMO start, with starts taken as the targets nearest to the depot.

\subsection{Anchor selection and solver-free expert search}

Two policies interact during training. The learning policy is updated at every step. A frozen reference policy generates search material and is replaced only when the learning policy obtains a strictly better exact mean on a fixed validation set of 1,000 instances. This gate ensures that the expert search starts from routes of validated quality.

An edge predictor maps the detached encoder output to a non-negative directed table of size $(J+1) \times (J+1)$. It is trained with a Huber loss towards tables built from exact solutions, and its loss does not reach the policy. Its only role is to rank the 200 reference routes and choose the anchor, so a poor prediction costs search quality and cannot affect validity.

The anchor table scores all sampled routes. We select 30 starts: the 20 routes with the lowest bound, and 10 further routes chosen for diversity from a pool of the 60 next-best routes within a relative gap of 5\%. From each start, directed 2-opt runs on the table for at most 30 moves and accepts only strict decreases of $U_{q}$. Local search on the table is complementary to LKH \citep{lin1973,helsgaun2000}: on large instances, a share of the starts is also passed to LKH as an asymmetric TSP on $C^{q}$. In every case the search minimises a surrogate that can only overestimate the true cost. The LKH time budget for all participating starts is deterministic, $1.5 \times 10^{-4} J^{2}$ seconds scaled by the LKH share, so that the search cost of a batch can be predicted in advance.

Proposition~\ref{prop:bound} gives the screening rule. If a route satisfies $U_{q}(\pi) < f(q)$, then $f(\pi) < f(q)$ holds without any solve, so the route is a certified improvement. Its exact cost is then computed once, and the route becomes the expert. If no route clears the threshold, the anchor remains the expert. Because the bound is conservative, some truly better routes fail this test. With a small probability, the near-miss rule therefore solves the route with the smallest relative gap above the threshold and accepts it if it is exactly better. The probability adapts to training progress: for small instances it moves between 1\% and 5\% as the imitation accuracy and the share of policy routes that beat the expert rise. Each instance has a budget of one such verification, and the solve is charged whether or not the route is accepted. If the exact cost of a near-miss route exceeds its bound, the run stops, because that would falsify Proposition~\ref{prop:bound}.

\subsection{Learning from sparse exact feedback}

Each instance yields three kinds of cost information for the POMO routes $\pi^{1},\dots,\pi^{P}$ of the learning policy. The anchor table and the expert table bound every route, and the code takes the route-wise minimum of the two. The route with the lowest bound is then solved exactly. Writing $c^{p}$ for the exact cost of that route and for the bound of every other route, the dense advantage is
\begin{equation}
a^{p} = \frac{\bar{c} - c^{p}}{\max\{\sigma(c),\, 10^{-6}\}},
\end{equation}
where $\bar{c}$ and $\sigma(c)$ are the mean and standard deviation over the POMO routes of the instance. Standardising within each instance prevents hard instances from dominating the gradient. A second term compares the exactly solved policy route $\pi^{\mathrm{sel}}$ with the verified expert $\pi^{E}$,
\begin{equation}
a^{\mathrm{ex}} = \operatorname{clip}\!\left(\frac{f(\pi^{E}) - f(\pi^{\mathrm{sel}})}{|f(\pi^{E})|},\, -1,\, 1\right).
\end{equation}
The reinforcement loss combines the two terms using REINFORCE \citep{williams1992}:
\begin{equation}
\mathcal{L}_{\mathrm{RL}} = -\frac{1}{BP}\sum_{b,p} a^{p}_{b} \log p_{\theta}(\pi^{p}_{b}) - \frac{w}{B} \sum_{b} a^{\mathrm{ex}}_{b} \log p_{\theta}(\pi^{\mathrm{sel}}_{b}),
\end{equation}
where $B$ is the batch size. The weight $w$ rises linearly from 0 to 1 as the imitation accuracy of the previous epoch moves from 0.50 to 0.85, so the comparison with the expert begins only when the policy can produce competitive routes.

A separate optimiser step performs teacher-forced cross-entropy on the expert sequences. Each batch mixes current experts with a replay buffer of 50,000 experts, which supplies 30\% of the batch. Only verified experts enter this update, so an error in the bound can waste a solve but can never become a label. Imitation runs for five inner epochs at each update on batches of 1,000 sequences, split into micro-batches of 128. The two updates use a common network and take separate optimiser steps.

\subsection{Progressive search}
\label{sec:recipe}

Search effort is scheduled by instance size and training stage, because a strong expert helps large instances and hurts small ones. Earlier runs showed that small instances reach their best validation value early and then deteriorate, a pattern we traced to falling route diversity in the POMO output. Instances with $J < 60$ therefore keep the plain 2-opt search throughout. For $J \ge 60$, a ramp $r = \operatorname{clip}((e - 40)/60,\, 0,\, 1)$ at epoch $e$ raises the fraction of LKH starts to $r/2$ and raises the near-miss probability towards 1. After epoch 60, large instances also receive two refresh rounds: the most recently solved route supplies a new table, which is tight around that route, and the search continues on it. Each round solves at most one additional route.

The quantity of imitation labels is normalised at the same time. An improved expert is repeated in the imitation batch so that the expected number of samples per instance stays near 1.44, with the repeat weight clipped to $[1, 5]$. Without this step, the imitation pressure would grow as the search improves.

The values of these settings came from earlier development runs. We compared the expert operators on a frozen policy and found that swap, 2-opt and LKH were complementary. Those runs also showed that the benefit of a stronger expert concentrates at $J \ge 60$, which set the size threshold. The ramp over epochs 40 to 100 was chosen so that the policy first learns from the cheapest search and receives the stronger expert only once it has stabilised. We did not run a formal search over these values, and Section~\ref{sec:discussion} returns to this point.

\subsection{Solver budget}

Per instance, the base loop uses one anchor solve, one target solve and at most one expert solve, and refresh rounds add at most two solves for large instances. The training run made 3,839,970 calls for 1,330,000 instances, or 2.89 calls per instance. Validation added 266,000 calls. By contrast, POMO with an exact reward makes one call per route per instance, so its cost grows with $J$. The composition of these calls is reported in Section~\ref{sec:experiments}, and the training loop is otherwise unconstrained by the solver: local search on the tables took 67.6 hours of summed worker time, against 86.7 hours of Gurobi time.

\section{Computational experiments}
\label{sec:experiments}

\subsection{Setup}

\paragraph{Baseline.} G200 is POMO trained for 200 epochs with an exact Gurobi reward on every sampled route, on 100,000 random instances per epoch. It uses the same instance generator and POMO start rule as the proposed method, a smaller network (6 encoder layers, embedding dimension 128, 8 heads, 1.27 million parameters, against 2.44 million for E133), and a batch size of 512. The two models therefore differ in architecture as well as in training signal, and the comparison in this paper is between complete training recipes. It was not retrained for this paper. An independent retraining under the same recipe is reported in Section~\ref{sec:ablation}.

\paragraph{Proposed model.} Training ran for 200 epochs of 10,000 instances each, with a batch size of 64, Adam with learning rate $10^{-4}$ and weight decay $10^{-6}$, and a cosine schedule. We report the checkpoint at epoch 133 (E133), selected by mean exact cost on a validation set of 2,000 instances, with 400 at each of $J \in \{20,40,60,80,100\}$. The selection used no test data. One fixed-order SOCP is solved by Gurobi with one thread per call. Experiments ran on a workstation with an Intel i9-13900KF processor (32 hardware threads), 31\,GB of memory and an NVIDIA RTX 4090 GPU. Over the selection range, validation means were close: E133 reached 17.703 and E150 reached 17.705.

\paragraph{Data.} Training instances have $J$ drawn uniformly from $\{20,\dots,100\}$ and targets drawn uniformly from the unit square. Speeds and endurance are sampled uniformly from ranges derived from the literature benchmark, with $v_{c} = 0.2$, $v_{d} \in [0.4, 1]$ and $E \in [0.05, 1]$ in normalised units. We evaluate on three sets. Test500 contains 100 new instances at each of $J \in \{20,40,60,80,100\}$. The generalisation set contains 100 new instances at each of $J \in \{150,200,300,500\}$. Real529 is the literature benchmark of \citet{gambella2018}, \citet{erdogan2021} and \citet{poikonen2020}, as compiled by \citet{li2025}, with 72, 132 and 325 instances. Its sizes include $J$ from 10 to 19 and $J = 200$, which lie outside the training range. Test500 serves as the primary evidence because it matches the training distribution and allows a clean comparison at each size. We inspected test500 once for an earlier checkpoint before selecting E133, and the selection itself relied on the validation set only. Real529 is not a held-out set, because earlier checkpoints of this project were evaluated on it during method development.

\paragraph{Test protocol.} Both models decode greedily with POMO, with $\min(J,100)$ starts on Real529 and $J$ starts on the synthetic sets. Under eight-view augmentation, each geometric view is decoded separately. Every distinct route is solved exactly and the best value is reported. Costs are therefore exact in all tables.

\paragraph{Metrics.} For model $A$ and baseline $B$ on instance $n$, the relative gap is $100\,(c^{A}_{n} - c^{B}_{n})/c^{B}_{n}$, so positive values mean that $A$ is worse. We report the mean over instances with a 95\% normal-approximation interval, together with the counts of wins, ties and losses at a tolerance of $10^{-3}$. We prefer the per-instance mean because a ratio of mean costs is dominated by the largest instances. On Real529, for example, the ratio of means is +0.12\%, whereas the per-instance mean is +0.21\%.

\subsection{Results on synthetic instances}

Table~\ref{tab:test500} reports test500. Without augmentation, E133 is behind G200 at every size. The gap is +1.08\% at $J = 20$, rises to about +3.4\% at $J = 60$ and $J = 80$, and is +3.05\% at $J = 100$; the mean over all sizes is +2.60\%. G200 wins 407 of 500 instances. This is the largest gap in our evaluation, and it marks the main quality cost of the proposed training scheme. The gap does not grow steadily with size. It is smallest at $J = 20$ and stabilises between $J = 60$ and $J = 100$, which is the range in which the progressive search is active. The mean costs at $J = 20$ differ by about 0.12 units, and ties are common there, whereas at $J \ge 60$ nearly all instances are lost by a modest margin. The picture is one of a consistent shortfall and not of occasional bad instances: E133 wins only 10 to 14 instances out of 100 at the three largest sizes.

\begin{table}[t]
\centering
\caption{Test500 without augmentation. Gap is E133 relative to G200 (positive means E133 is worse).}
\label{tab:test500}
\small
\begin{tabular}{rrrrrr}
\toprule
$J$ & E133 & G200 & Gap (\%) & 95\% interval & Win/Tie/Loss \\
\midrule
20 & 12.594 & 12.472 & +1.08 & [+0.72, +1.43] & 15/22/63 \\
40 & 15.923 & 15.633 & +2.00 & [+1.55, +2.45] & 18/1/81 \\
60 & 18.502 & 17.946 & +3.39 & [+2.77, +4.01] & 13/0/87 \\
80 & 20.079 & 19.453 & +3.47 & [+2.98, +3.96] & 10/0/90 \\
100 & 22.040 & 21.440 & +3.05 & [+2.43, +3.68] & 14/0/86 \\
All & 17.828 & 17.389 & +2.60 & [+2.35, +2.84] & 70/23/407 \\
\bottomrule
\end{tabular}
\end{table}

Eight-view augmentation narrows the gap in the runs available so far. \placeholder{E133 with eight views on test500 has not been run. Available result: the earlier checkpoint E150 has a mean gap of +1.07\% to G200 under eight views, which must not be reported as E133. Replace this paragraph with the E133 result.}

\subsection{Generalisation to larger instances}

Table~\ref{tab:gen} reports instances with 150 to 500 targets, which are beyond the training range. Under eight views, the mean gap to G200 is +0.72\%, and at $J = 300$ the interval includes zero. These gaps are smaller than the unaugmented test500 gaps, but the two protocols differ, so the tables are not directly comparable. Without augmentation, the mean gap on the same instances is +1.63\%. Augmentation helps both models, and it helps E133 more: the unaugmented mean is 36.801 for E133 and 36.336 for G200, and eight views reduce these to 35.827 and 35.644. Neither model was trained at these sizes. The results thus show that learning from the table did not cost additional transfer ability relative to learning from exact rewards.

\begin{table}[t]
\centering
\caption{Generalisation set with eight views. Gap is E133 relative to G200.}
\label{tab:gen}
\small
\begin{tabular}{rrrrrr}
\toprule
$J$ & E133 & G200 & Gap (\%) & 95\% interval & Win/Tie/Loss \\
\midrule
150 & 26.003 & 25.795 & +0.96 & [+0.62, +1.31] & 23/0/77 \\
200 & 29.077 & 28.868 & +0.97 & [+0.56, +1.37] & 30/0/70 \\
300 & 36.945 & 36.875 & +0.38 & [$-$0.03, +0.78] & 42/0/58 \\
500 & 51.283 & 51.040 & +0.56 & [+0.23, +0.90] & 40/0/60 \\
All & 35.827 & 35.644 & +0.72 & [+0.53, +0.91] & 135/0/265 \\
\bottomrule
\end{tabular}
\end{table}

\subsection{Results on the literature benchmark}

Table~\ref{tab:real} reports Real529 by size band. The mean gap to G200 is +0.21\%, with an interval that excludes zero. Gaps are below +0.53\% in every band up to $J = 100$. At $J = 200$, E133 is 1.25\% better than G200, with 21 wins out of 25. Ties are frequent at the smallest sizes, where 127 of 220 instances end at the same value, which suggests that both models often reach the same route. This holds for 25 instances from one source and lies outside the training range, so we treat it as an observation and not as evidence of general superiority: on the synthetic instances at the same size in Table~\ref{tab:gen}, E133 is behind.

\begin{table}[t]
\centering
\caption{Real529 with eight views, by size band. Gap is E133 relative to G200.}
\label{tab:real}
\small
\begin{tabular}{lrrrrrr}
\toprule
Band & $n$ & E133 & G200 & Gap (\%) & 95\% interval & Win/Tie/Loss \\
\midrule
$J = 10$--19 & 220 & 255.570 & 255.086 & +0.19 & [+0.12, +0.26] & 23/127/70 \\
$J = 20$--50 & 234 & 342.841 & 341.847 & +0.33 & [+0.21, +0.45] & 60/33/141 \\
$J = 100$ & 50 & 431.749 & 429.727 & +0.52 & [+0.13, +0.92] & 19/0/31 \\
$J = 200$ & 25 & 720.934 & 730.123 & $-$1.25 & [$-$1.78, $-$0.72] & 21/0/4 \\
All & 529 & 332.819 & 332.421 & +0.21 & [+0.13, +0.29] & 123/160/246 \\
\bottomrule
\end{tabular}
\end{table}

Table~\ref{tab:meta} compares both models with best known solutions and with three metaheuristics that search over visiting orders and score every order with the same exact SOCP. The reference value of each instance is the best upper bound over all methods in \citet{li2025}. \placeholder{Merge the best values of Vieira et al. (2026) into the reference.} Each metaheuristic was given a multiple of the single-core time that E133 needs on that instance. At three times that budget, iterated local search (ILS) and variable neighbourhood search (VNS) are worse than E133. At ten times the budget they are better than both learned policies, at +0.73\% and +0.78\% against +1.35\% for E133 and +1.14\% for G200. The difference is concentrated in the clustered instances of \citet{poikonen2020}, where VNS reaches $-$0.03\% and E133 is at +2.10\%. Clustered instances differ from the uniform distribution used for training, which is consistent with the weakness of both learned policies there. Simulated annealing (SA) remains behind at all budgets, with the gap at +3.14\% at ten times the budget. On average, E133 and G200 need 0.78\,s of wall time per instance, or 10.3\,s on a single core, on the same hardware. The time grows with size, from under 0.1\,s of wall time at $J = 10$ to about 5.6\,s at $J = 200$, with 16 parallel solver workers. Relative to the best values in \citet{li2025}, both models also find new best solutions on a number of Poikonen instances, for which the reference values come from a heuristic. We have recomputed one such case geometrically and found it feasible under our model. \placeholder{Report the counts of new best solutions once the reference includes Vieira et al. (2026).}

\begin{table}[t]
\centering
\caption{Real529 gap to the best known solution (BKS), $100\,(\mathrm{UB}-\mathrm{BKS})/\mathrm{UB}$, where UB is the upper bound found by the method, by subset. Budgets are multiples of the single-core time of E133.}
\label{tab:meta}
\small
\resizebox{\textwidth}{!}{\begin{tabular}{lrrrrrr}
\toprule
Method & Gam72 & \begin{tabular}[c]{@{}c@{}}Erd132\\ 16--20\end{tabular} & \begin{tabular}[c]{@{}c@{}}Erd132\\ 25--50\end{tabular} & \begin{tabular}[c]{@{}c@{}}Poi325\\ Regular\end{tabular} & \begin{tabular}[c]{@{}c@{}}Poi325\\ Cluster\end{tabular} & All \\
\midrule
E133 & +0.19 & +0.24 & +0.89 & +1.75 & +2.10 & +1.35 \\
G200 & +0.01 & +0.08 & +0.80 & +1.59 & +1.76 & +1.14 \\
ILS, $3\times$ & +0.71 & +0.80 & +3.08 & +2.79 & +0.36 & +1.63 \\
VNS, $3\times$ & +0.69 & +0.76 & +3.14 & +2.84 & +0.36 & +1.65 \\
ILS, $10\times$ & +0.33 & +0.27 & +1.90 & +1.22 & $-$0.02 & +0.73 \\
VNS, $10\times$ & +0.20 & +0.34 & +1.89 & +1.40 & $-$0.03 & +0.78 \\
SA, $10\times$ & +0.53 & +1.27 & +5.15 & +5.16 & +1.84 & +3.14 \\
\bottomrule
\end{tabular}}
\end{table}

\subsection{Comparison with a general-purpose exact solver}

To set the scale of the problem, we also solved the full mixed-integer SOCP (Model+ of \citet{li2025}) with Gurobi, with a time limit of 3,600\,s per instance and one thread. \placeholder{Report status, upper bound, lower bound, gap and time by $J$ for Real529 and test500. The Real529 run has finished 489 of 529 instances and the test500 run has not started.}

\subsection{Training efficiency}

Table~\ref{tab:budget} compares training budgets. E133 used 15 times fewer training instances and 78 times fewer Gurobi calls than G200, or 73 times fewer when in-training validation calls are included. The test-time cost is unchanged in our measurements because it is dominated by exact solves of the POMO routes: on Real529, decoding takes 0.12\,s per instance on the GPU and Gurobi takes about 10.2\,s of single-core time for either model. CPU time for G200 is an estimate, obtained by multiplying its call count by a mean solve time measured at about $J = 60$. Wall time for G200 comes from a different machine with 128 reward workers, so the ratio of wall times shows the order of magnitude only. Solver calls and instances are exact counts and form the primary comparison.

Table~\ref{tab:calls} breaks down the 3.84 million calls of E133. Anchor and target solves contribute one call per instance each. Expert verification, including near-miss checks, contributes 0.37 calls per instance. Refresh rounds add 0.52 calls per instance on average across all sizes, and more than one per instance at the end of the run, when more than half of the batches have $J \ge 60$ and refresh is active. A share of 23.7\% of refresh solves produced a better expert, so the refresh budget has a measurable return but also a large waste. Near-miss checks were attempted 364,974 times.

\begin{table}[t]
\centering
\caption{Exact solver calls of E133, epochs 1 to 133, by purpose.}
\label{tab:calls}
\small
\begin{tabular}{lrr}
\toprule
Purpose & Calls & Per instance \\
\midrule
Anchor solve & 1,330,000 & 1.00 \\
Target reward solve & 1,330,000 & 1.00 \\
Expert and near-miss verification & 493,128 & 0.37 \\
Refresh rounds & 686,842 & 0.52 \\
Total & 3,839,970 & 2.89 \\
\bottomrule
\end{tabular}
\end{table}

Table~\ref{tab:dynamics} shows how the loop behaves during training. In the first epoch, almost any search improves on an untrained policy. By epoch 40 the plain 2-opt search improves only 9\% of the anchors, by 0.38\% on average. The ramp then lifts the share of improved experts above one third, which is the intended effect of the progressive search. Over the same period, the lowest-bound policy route has a bound that falls from 92\% to 34\% above its exact cost, which is consistent with policy routes moving closer to the anchors as training proceeds. The bound remains loose at epoch 133, yet training continues to make progress; Section~\ref{sec:discussion} discusses why.

\begin{table}[t]
\centering
\caption{Training dynamics of the proposed method. Improved share is the fraction of instances whose expert beats the anchor. Bound excess compares the lowest bound among policy routes with the exact cost of that route.}
\label{tab:dynamics}
\small
\begin{tabular}{rrrrr}
\toprule
Epoch & Anchor cost & Expert cost & Improved share (\%) & Bound excess (\%) \\
\midrule
20 & 19.986 & 19.779 & 12.1 & 92.4 \\
40 & 19.205 & 19.096 & 9.0 & 70.7 \\
60 & 19.146 & 18.798 & 24.4 & 57.5 \\
80 & 18.993 & 18.640 & 33.0 & 57.9 \\
100 & 18.850 & 18.603 & 34.9 & 44.2 \\
133 & 18.522 & 18.288 & 37.9 & 33.6 \\
\bottomrule
\end{tabular}
\end{table}

\begin{table}[t]
\centering
\caption{Training budgets of G200 and E133.}
\label{tab:budget}
\small
\resizebox{\textwidth}{!}{\begin{tabular}{lrrr}
\toprule
 & G200 & E133 & Ratio \\
\midrule
Training instances & 20,000,000 & 1,330,000 & 15 \\
Gurobi calls in training & 299,437,760 & 3,839,970 & 78 \\
Gurobi CPU (core-hours) & about 2,580 (estimated) & 86.7 & about 30 \\
Wall time (hours) & about 215.5 (128 workers) & 19.2 (32 workers) & about 11 \\
Test time per Real529 instance (s) & 0.78 & 0.78 & 1 \\
\bottomrule
\end{tabular}}
\end{table}

\placeholder{Add the 24-hour training curve: validation cost against wall time for POMO with exact rewards and for the proposed method on the same hardware. The POMO run is queued; confirm the plan for the proposed-method curve.}

\subsection{Ablation study}
\label{sec:ablation}

The ablations answer three questions about the design. Variants B and D train from scratch to one million instances (epoch 100) and are compared with the proposed method at the same point on the validation set of 2,000 instances; the results are then reported on test500 by size.
\begin{itemize}
\item Group A retrains POMO with exact rewards from scratch under the G200 recipe, on the same hardware, for 24 hours. Its checkpoints at the instance and solver budgets of E133 measure the quality attainable at a matching budget.
\item Group B replaces table-based search with exact evaluation of candidate moves, under a budget of 100 exact solves per instance, to measure what the table saves in training time.
\item Group D removes expert diversity: one greedy route receives the full local-search budget of the 30 starts, to separate the effect of search diversity from that of search effort.
\end{itemize}
Group C is the proposed method at the same point in training. \placeholder{Report quality, wall time and exact calls for groups A, B, C and D. All are queued.}

\section{Managerial insights and discussion}
\label{sec:discussion}

\subsection{When the approach is worthwhile}

The method changes when computing effort is spent. A metaheuristic spends it on every instance; a trained policy spends it once, and each later instance takes seconds. For a fleet operator that replans repeatedly, this favours a trained policy when the number of decisions is large or the replanning interval is short. The relevant cost for the operator is the training budget, and our results reduce the solver effort by nearly two orders of magnitude. An operator can therefore retrain for a new vessel speed, drone endurance or service region within a day on a single workstation, where the baseline recipe needed more than a week on a 128-worker setup. The method needs no training labels from an external solver run and no precomputed instance library, because it generates its own experts during training.

The cost of this saving depends on the use case. Where an operator has minutes per instance, the iterated local search results in Table~\ref{tab:meta} suggest that a metaheuristic will give better plans, particularly when targets are clustered. Where the plan must be available in seconds, the learned policy is the better option at the time budgets we tested: ILS and VNS at three times the single-core time of E133 are behind it. The two can also be combined, since a route from the policy can start a metaheuristic. We did not test that combination.

\subsection{Sensitivity to operating parameters}

\placeholder{Sensitivity analysis on speed ratio, drone endurance and target distribution has not been run. Add results and practical implications here, or remove this subsection.}

\subsection{Why a loose bound suffices}

The bound is far from tight on routes that differ from the anchor: Table~\ref{tab:dynamics} shows an excess of 33.6\% at epoch 133. In a continuation experiment from E133, we took the minimum over all tables available for each route, which used about two additional tables per instance on the affected batches and tightened the bound by only 0.32\%. It left validation quality unchanged. A much tighter bound is thus not easy to obtain from nearby solutions, and on the evidence so far it is not the limiting factor. Two properties of the design explain why a loose bound can work. The bound is exact at the anchor and accurate nearby, which is where local search operates. The labels that reach imitation learning are exact, so the looseness affects only the ranking of candidates and the variance of the reward. The relevant property of the table is that it is cheap and valid, not that it is tight. Group B of the ablation will quantify this claim. This observation has a wider implication for the design of learning methods with expensive objectives: a signal with a guaranteed direction of error and a verification step can substitute for an accurate one. It is not an argument against tighter bounds. It says only that the marginal value of tightness appears low in the region we tested.

\subsection{Limitations}

Five limitations apply. First, E133 remains behind G200 on synthetic instances, by up to about 3.5\% without augmentation. The evidence points to route diversity. In diagnostics on a later checkpoint from the same method (epoch 163), the overlap between POMO routes at $J = 100$ was 0.62, against 0.53 for G200. Attempts to recover diversity by changing the reward bound or the imitation weight had little effect. Second, the results rest on one training seed. \placeholder{Add two further seeds.} Third, the comparison on time depends on hardware and implementation, which is why we treat solver calls as the primary measure. Fourth, the settings of the progressive search were set by development runs and not by a formal search, so the reported performance may not be the best attainable for the method. Fifth, the bound depends on the separable structure of the CVTSP. It extends to a variant only if sortie constraints remain independent of the order, which should hold for several sorties of fixed composition but fails if, for instance, target time windows couple the sorties.

\section{Conclusion}
\label{sec:conclusion}

This paper asked whether a DRL policy for the CVTSP needs an exact solver reward for every route. It does not. One optimal fixed-order solution provides an upper bound on all visiting orders, computed by summing table entries, and this bound can supply most of the training signal. With 2.89 exact solves per training instance, the resulting policy uses 15 times fewer instances and 78 times fewer solver calls than POMO trained on exact rewards, at the same test time. Its quality is close to that baseline on the literature benchmark and on large instances, with mean gaps of +0.21\% and +0.72\%, and it remains behind on synthetic instances within the training range, with a mean gap of +2.60\% without augmentation. Metaheuristics given ten times the computing time of the policy on each instance obtain better solutions on the literature benchmark, so the method suits settings where plans are needed quickly and where retraining has to be affordable.

The idea carries over to other routing problems in which a route cost comes from a continuous subproblem with a separable structure, because all it requires is that a solution for one route remain feasible for others. Future work should target the loss of route diversity, which appears to be the main source of the remaining gap, and should test whether the table can provide fast warm starts for metaheuristics. Extending the approach to several carriers, as in the vessel routing model of \citet{ren2026}, would test how far the separable structure reaches.

\section*{CRediT authorship contribution statement}
\placeholder{To be completed by the authors.}

\section*{Declaration of competing interest}
\placeholder{To be completed by the authors.}

\section*{Funding}
\placeholder{To be completed by the authors.}

\section*{Data availability}
\placeholder{State where the code, generated instances and per-instance results are available. Real529 instances are from the cited sources.}

\section*{Ethics statement}
This study used computational experiments on synthetic and published benchmark instances and involved no human participants or personal data.

\section*{Declaration of generative AI use}
\placeholder{The authors should confirm this statement.} During the preparation of this work the authors used Claude (Anthropic) to assist with literature search and manuscript drafting. The authors reviewed and edited the content and take full responsibility for the publication.

\bibliographystyle{elsarticle-harv}
\bibliography{refs}

\end{document}
